\section{Introduction}

If artificial intelligence displaces a substantial share of UK employment, who bears the resulting loss: displaced workers, lower-income households or the Exchequer? The answer depends not only on the size of the shock but also on how it is distributed. The UK's services-intensive economy has substantial employment in professional, financial, administrative and managerial occupations that score highly on common measures of AI exposure \citep{felten2021, eloundou2023, dsit2023, henseke2025}. Its tax-benefit system combines means-tested support with substantial taxation of labour income, so the household and fiscal effects depend on which workers lose earnings.

Exposure indices do not, by themselves, establish that displacement will be concentrated towards the top of the earnings distribution. Evidence on AI's labour-market incidence remains unsettled. Exposure-based scenarios assign greater risk to highly paid cognitive work \citep{felten2021, eloundou2023}; an alternative account emphasises compression of expertise premia, with less experienced workers gaining access to tasks previously performed by more highly paid specialists \citep{autor2024}; and several early empirical studies report weaker employment outcomes for junior workers in exposed occupations or firms \citep{kleinteeselink2025, hosseini2026, brynjolfsson2025}. These are distinct hypotheses about incidence rather than refinements of a single estimate. They imply different household losses, benefit responses and reductions in tax receipts.

I treat incidence as an explicit dimension of the scenario design. Using PolicyEngine UK and processed Family Resources Survey 2024--25 microdata, I assign occupation-level AI exposure to workers and compare five displacement allocations: exposure-proportional, junior-concentrated, expertise-compression, uniform, and a Klein-anchored top-loaded stress test. The simulation adapts the Irish tax-benefit framework of \citet{doorley2026}, including its 7 per cent employment-loss and 2.6 per cent survivor-wage conventions, while extending the design by varying incidence explicitly. These values convert quantities discussed by \citet{briggs2023}; they are not estimates from that study. I also simulate 66 combinations of displacement rates and wage uplifts, together with a separate wage-adjustment family in which author-imposed occupation-level cuts replace job losses. The code and aggregate artefacts are public, but exact regeneration requires licensed FRS microdata. Plain survey weights and the absence of the official SPI top-income adjustment limit representativeness at the top of the distribution.

Two findings organise the paper. First, measured inequality rises in every displacement calibration examined: across all 66 shock-size cells, the Gini change ranges from \genGridGiniMinPp{} to \genGridGiniMaxPp{} percentage points; across the five central incidence families, the 20-draw means range from \genIncFamilyGiniMeanMinPp{} to \genIncFamilyGiniMeanMaxPp{} points. This is a conditional result of the simulated displacement, survivor-wage and capital-income channels, not a claim about every possible AI adjustment. In the exposure-proportional case, selected higher-income households lose earnings while surviving workers receive wage gains and capital-income gains remain concentrated near the top. Lower tax liabilities cushion some losses but do not prevent the simulated disposable-income distribution from widening.

Second, incidence changes the balance between fiscal and poverty effects. Holding the weighted displacement quota and central parameters fixed, the 20-draw mean Exchequer cost ranges from \pounds \genIncUniformCostMeanBn~billion under uniform incidence to \pounds \genIncCompressionCostMeanBn~billion under expertise compression and \pounds \genIncKleinCostMeanBn~billion under the Klein-anchored stress test. Across the four stylised families, absolute before-housing-costs poverty rises by \genIncFamilyPovBhcMeanMinPp{} to \genIncCompressionPovBhcMeanPp{} percentage points; because the families share assignment seeds, paired draw-level comparisons show the expertise-compression poverty effect reliably exceeding the exposure-proportional one (by $+0.13$ points, paired 95 per cent interval $[+0.04, +0.23]$), while the uniform--exposure poverty difference ($+0.02$ points, $[-0.08, +0.11]$) and the junior--exposure fiscal-cost difference are statistically indistinguishable. % paired differences recomputed from results/robustness/incidence_draws_five.csv (20 shared seeds, paired normal CI)
The adjustment margin also matters. In the paired seed-0 comparison, moving the same 7.77 per cent gross earnings loss from wage cuts to job loss lowers the simulated fiscal cost from \pounds \genWageMarginCentralCostBn~billion to \pounds \genCentralCostBn~billion, while raising poverty from 0.10 to 1.81 points and the Gini change from $-0.44$ to $+1.04$ points. % results/wage_margin_central.json vs results/central.json
These comparisons quantify the consequences of the specified incidence and adjustment rules; they do not identify which pattern will occur.



The remainder of the paper reviews the literature, sets out the scenario and microsimulation methods, reports the distributional and fiscal results, examines three stylised policy responses, and concludes with limitations and implications.
